\documentclass[10pt]{article}

\usepackage[letterpaper,margin=0.67in]{geometry}
\usepackage{amsmath,amssymb}
\usepackage{booktabs}
\usepackage{microtype}
\usepackage[hidelinks]{hyperref}
\usepackage{enumitem}

\setlength{\parindent}{0pt}
\setlength{\parskip}{3pt}
\setlist{nosep,leftmargin=1.25em}
\setcounter{secnumdepth}{2}
\renewcommand{\baselinestretch}{0.96}
\newcommand{\R}{\mathbb{R}}

\title{\vspace{-1.7em}\textbf{Tubelet-ViT LeJEPA for Volumetric Brain MRI Pretraining}}
\author{MRIxFields2026 Group Project}
\date{}

\begin{document}
\maketitle
\vspace{-2.2em}

\begin{abstract}
This report presents a volumetric adaptation of LeJEPA for representation
learning from cross-field brain MRI. Instead of encoding axial slices
independently, the method constructs related views of a 16-slice anatomical
slab and tokenizes them with a three-dimensional tubelet embedding. A single
ViT-B encoder maps all views to projected CLS representations trained with
LeJEPA's global-center invariance objective and Sketched Isotropic Gaussian
Regularization (SIGReg). The resulting artifact is an encoder, not a complete
image translator: its dense three-dimensional patch tokens are intended to
initialize a later conditional MRI synthesis model. The design therefore
preserves the simplicity of LeJEPA while explicitly monitoring the unresolved
gap between pooled self-supervision and dense anatomical transfer.
\end{abstract}

\section{Motivation and Scope}

MRIxFields2026 studies synthesis between five magnetic field strengths and
three structural MRI modalities \cite{mrixfields}. The preceding project
experiment applied LeJEPA to individual axial slices. Although this provided a
controlled two-dimensional baseline, a slice-wise encoder cannot directly use
through-plane continuity, and its patch features were only indirectly trained
through a pooled CLS loss. The present method changes the input representation
and encoder while retaining the established LeJEPA objective.

LeJEPA combines agreement between related views with a distributional
regularizer that prevents complete and dimensional collapse without a target
network, exponential moving average (EMA), stop-gradient, or contrastive
negatives \cite{lejepa}. This simplicity is the main reason for retaining it.
Tubelet tokenization does not turn the method into V-JEPA: V-JEPA and EchoJEPA
use masked latent prediction with a context encoder, predictor, and EMA target
encoder \cite{vjepa,echojepa}, whereas the selected method remains an
unmasked, view-based LeJEPA model.

The scope is volumetric self-supervised pretraining and encoder export. Field
and modality labels are used for balanced sampling and reporting, but no
source-field, target-field, or modality conditioning is introduced during
self-supervision. Cross-field control belongs to the downstream translator,
where a target field is well defined.

\section{Volumetric Multi-View Formulation}

Let a base sample be a contiguous axial slab
\begin{equation}
  s_n=[x_z,x_{z+1},\ldots,x_{z+15}]
  \in [0,1]^{1\times16\times H\times W}.
\end{equation}
All related views retain the same subject, field strength, modality, and
anatomical neighborhood. Two global views have shape
$1\times16\times224\times224$. Six local views use contiguous 8-slice
sub-slabs and have shape $1\times8\times96\times96$. Their depth offsets are
distinct, and each local sub-slab contains the common anchor anatomy needed by
the intended center-slice consumer.

LeJEPA supplies the two-global and multi-local alignment structure
\cite{lejepa}. EchoJEPA provides a medical-video precedent for 16-frame inputs,
tubelet depth 2, patch size 16, global crop scale $[0.5,1.0]$, and aspect ratio
$[0.9,1.1]$ chosen to avoid distorting clinically meaningful geometry
\cite{echojepa}. Our local scale $[0.1,0.3]$, foreground-aware overlap, shared
anchor coordinate, and anchor-containing depth windows are MRI-specific
project choices rather than claims about those papers.

Each view samples one crop and affine transform and applies it coherently to
every slice. Small rotations, translations, mild blur, and mild Gaussian noise
are allowed. Natural-image color jitter, grayscale conversion, solarization,
large rotations, and aggressive contrast changes are excluded because field
strength and modality appearance are part of the later translation signal.

\newpage

\section{Tubelet-ViT LeJEPA}

\subsection{Tubelet encoder}

The encoder begins with a non-overlapping three-dimensional convolution:
\begin{equation}
  \mathcal{P}(X)=\operatorname{Conv3D}_{k=s=(2,16,16)}(X),
  \qquad C_{\mathrm{out}}=768.
\end{equation}
For a global view this gives
\begin{equation}
  X^{(g)}\in\R^{B\times1\times16\times224\times224}
  \longrightarrow
  T^{(g)}\in\R^{B\times8\times14\times14\times768},
\end{equation}
or 1,568 patch tokens per view. A local view produces a
$4\times6\times6$ grid, or 144 tokens. AV-JEPA is the closest direct
architectural precedent: it tokenizes 16 video frames using a
$2\times16\times16$ Conv3D stem and feeds the resulting 1,568 tokens to a
shared ViT-B \cite{avjepa}. VideoMAE and V-JEPA provide the broader
space-time tubelet lineage \cite{videomae,vjepa}. Our adaptation changes RGB
video to one-channel MRI and interprets the temporal dimension as ordered
anatomical depth.

A learned CLS token is prepended before a standard bidirectional ViT-B with 12
blocks, hidden width 768, 12 attention heads, and MLP ratio 4
\cite{vit}. The global $8\times14\times14$ grid has a learned absolute
three-dimensional positional table. Patch positions are trilinearly
interpolated for local and downstream grids while the CLS position remains
separate. The same tubelet stem, transformer, and final LayerNorm process every
view.

For view $v$ of sample $n$, the training representation is
\begin{equation}
  h_{n,v}=f_{\theta}(v_{n,v})_{\mathrm{CLS}},\qquad
  z_{n,v}=g_{\phi}(h_{n,v})\in\R^{128},
\end{equation}
where $g_{\phi}$ is a shared BatchNorm MLP with dimensions
$768\rightarrow2048\rightarrow2048\rightarrow128$. Only the projected CLS
enters the SSL loss. Dense post-LayerNorm patch tokens remain available from
$f_{\theta}$ for transfer.

\subsection{Global-center invariance and SIGReg}

Let $V_g=2$ be the number of global views and $V=8$ the total number of views.
The target for sample $n$ is the mean of its global projections:
\begin{equation}
  \mu_n=\frac{1}{V_g}\sum_{v=1}^{V_g}z_{n,v}.
\end{equation}
All views predict this center through
\begin{equation}
  \mathcal{L}_{\mathrm{inv}}
  =\frac{1}{BV}\sum_{n=1}^{B}\sum_{v=1}^{V}
   \left\|z_{n,v}-\mu_n\right\|_2^2.
\end{equation}
There is no detach operation on the center or either encoder path.

For each view, SIGReg projects the physical batch onto random unit directions
$a\in\mathcal{A}\subset\mathbb{S}^{D-1}$ and applies the Epps-Pulley
Gaussianity statistic $T_{\mathrm{EP}}$:
\begin{equation}
  \mathcal{L}_{\mathrm{SIGReg}}
  =\frac{1}{V|\mathcal{A}|}
   \sum_{v=1}^{V}\sum_{a\in\mathcal{A}}
   T_{\mathrm{EP}}\!\left(\{a^{\top}z_{n,v}\}_{n=1}^{B}\right).
\end{equation}
The regularizer is computed separately across the physical batch of each view;
correlated views are not flattened into the sample axis. The complete objective
is \cite{lejepa}
\begin{equation}
  \mathcal{L}_{\mathrm{LeJEPA}}
  =(1-\lambda)\mathcal{L}_{\mathrm{inv}}
  +\lambda\mathcal{L}_{\mathrm{SIGReg}},\qquad \lambda=0.05.
\end{equation}
The implementation retains 1,024 random directions and a 17-point quadrature
over the Epps-Pulley statistic from the preceding MRI LeJEPA experiment.

\section{Data and Training Formulation}

Self-supervised optimization uses only the retrospective cohort. It covers all
15 field-strength and modality domains. Prospective paired subjects are
excluded so that later subject-held-out translation experiments remain unseen
to both SSL and supervised fine-tuning. Sampling is hierarchical: first a
domain is chosen uniformly, then a volume within that domain, and finally a
valid 16-slice interval. This prevents domains or scans with more eligible
windows from dominating the representation objective.

The organizer-provided finite \texttt{float32} intensities remain in $[0,1]$
through view construction. After mild blur or noise they are clamped and mapped
to the encoder range by
\begin{equation}
  x_{\mathrm{enc}}=2x-1.
\end{equation}
No per-volume standardization, histogram matching, or field-specific
normalization is used because those operations could erase acquisition
appearance needed downstream. The encoder and projector are initialized from
scratch and optimized with AdamW in BF16.

\section{Representation Validation and Intended Transfer}

The objective acts on projected CLS embeddings, but the translator will consume
dense tokens. A CLS loss does backpropagate through attention and the tubelet
stem; it does not require each spatial coordinate to preserve local anatomy.
The SSL Cookbook likewise cautions that representation rank is only a necessary
condition and that methods tuned for global recognition may lack localization
needed by dense tasks \cite{cookbook}.

Accordingly, validation separates three questions. First, finite losses,
gradient norms, feature standard deviation, singular spectra, and effective
rank detect optimization failure or collapse. For singular values $\sigma_k$
of a representation matrix, the effective rank is
\begin{equation}
  \operatorname{Rank}_{\mathrm{eff}}(Z)
  =\exp\!\left(-\sum_k p_k\log p_k\right),
  \qquad p_k=\frac{\sigma_k}{\sum_j\sigma_j}.
\end{equation}
Second, related-view retrieval and coordinate-matched dense-token
correspondence test whether the encoder retains subject anatomy beyond global
collapse avoidance. Third, the causal transfer test compares the same
downstream model initialized either from the exported encoder or randomly.

Only the exact tubelet encoder is exported. The projection MLP and SIGReg are
training-only components. A downstream consumer receives post-final-LayerNorm
dense features
\begin{equation}
  F\in\R^{B\times8\times h\times w\times768},
\end{equation}
adds source-field, target-field, and modality conditioning outside the SSL
encoder, and fine-tunes the complete encoder. This avoids repeating the main
confounds of the prior two-dimensional transfer experiment: an incompletely
matched feature path and a mostly frozen backbone.

\section{Discussion}

The design combines a paper-backed video tubelet architecture with the simpler
LeJEPA objective and MRI-specific view semantics. Its central hypothesis is
that neighboring axial slices provide useful volumetric context without
requiring masked reconstruction or an EMA teacher. Its central limitation is
equally explicit: healthy pooled embeddings do not prove dense anatomical
utility. Dense diagnostics and matched downstream initialization are therefore
part of the method's evaluation, not optional supporting evidence.

\vspace{-0.4em}
\footnotesize
\begin{thebibliography}{9}
\setlength{\itemsep}{-1pt}

\bibitem{mrixfields}
MRIxFields2026 Organizers.
\newblock \emph{MRIxFields2026: Cross-Field Brain MRI Synthesis Challenge}.
\newblock Official challenge website and repository, 2026.

\bibitem{lejepa}
R. Balestriero and Y. LeCun.
\newblock LeJEPA: Provable and scalable self-supervised learning without the heuristics.
\newblock \emph{arXiv:2511.08544}, 2025.

\bibitem{avjepa}
B. Robson, S. Mentu, W. Zhao, and A. Solin.
\newblock AV-JEPA: Extending LeJEPA to audio-visual self-supervised learning.
\newblock \emph{arXiv:2607.15295}, 2026.

\bibitem{echojepa}
A. Munim et al.
\newblock EchoJEPA: A latent predictive foundation model for echocardiography.
\newblock \emph{arXiv:2602.02603}, 2026.

\bibitem{vjepa}
A. Bardes et al.
\newblock Revisiting feature prediction for learning visual representations from video.
\newblock \emph{arXiv:2404.08471}, 2024.

\bibitem{videomae}
Z. Tong, Y. Song, J. Wang, and L. Wang.
\newblock VideoMAE: Masked autoencoders are data-efficient learners for self-supervised video pre-training.
\newblock \emph{Advances in Neural Information Processing Systems}, 2022.

\bibitem{vit}
A. Dosovitskiy et al.
\newblock An image is worth 16x16 words: Transformers for image recognition at scale.
\newblock \emph{International Conference on Learning Representations}, 2021.

\bibitem{cookbook}
R. Balestriero et al.
\newblock A cookbook of self-supervised learning.
\newblock \emph{arXiv:2304.12210}, 2023.

\end{thebibliography}

\end{document}
